\chapter{NIENTE DURA}\label{niente-dura}

\hfill\break

Mi svegliai sapendo di non essere pronto a rivederla.

Ero nella lussuosa casa estiva affittata da E.D. Lawton nel Berkshire
con il Sole che splendeva attraverso le tende di pizzo filigranato,
pensando: ``Basta stronzate.'' Ne ero stufo. Di tutte le stronzate
egocentriche degli ultimi otto anni, inclusa la storia con Candice
Boone, che aveva intuito molto prima di me le mie bugie illusorie. ``Sei
un po' fissato con questi Lawton'' mi aveva detto una volta. ``Dillo a
me!'' avrei voluto risponderle.

Onestamente non potevo affermare di essere ancora innamorato di Diane.
Anzi, il legame tra noi non era mai stato tanto inequivocabile. Eravamo
cresciuti dentro e fuori da quel legame, come viti che s'intrecciano
attorno a un pergolato, e nei momenti migliori il nostro era stato un
rapporto vero, un'emozione quasi spaventosa per la sua solennità e
maturità. Ecco perché ero così ansioso di nasconderla. Avrebbe
spaventato anche lei.

Mi scoprivo ancora a tenere con lei conversazioni immaginarie, di solito
a tarda sera, offrendo divagazioni al cielo senza stelle. Ero abbastanza
egoista da sentirne la mancanza ma anche sufficientemente sano di mente
da riconoscere che non eravamo mai stati davvero insieme. Ero pronto sul
serio a dimenticarla.

Ma non ero pronto a rivederla.

\separatorefiori

Jason era seduto in cucina mentre io mi preparavo la colazione. Aveva
bloccato la porta in modo che restasse aperta. Dolci bave di vento si
diffondevano nella casa. Stavo seriamente pensando di lanciare la
valigia nel bagagliaio della Hyundai e darmela a gambe. «Parlami di
questo NR» dissi invece.

«Non li leggi proprio i giornali?» mi chiese Jason. «Lì a Stony Brook li
tengono in isolamento, gli studenti di medicina?»

«Scusa, certo che so qualcosa dell'NR, per lo più quello che ho sentito
dai notiziari o colto dalle conversazioni nella mensa. So che NR sta per
``Nuovo Regno''. So che è un movimento cristiano (almeno nel nome)
ispirato dallo Spin e che è stato stigmatizzato dall'opinione pubblica e
dalle chiese conservatrici. So che attrae soprattutto i giovani e i
disillusi.» Un paio di ragazzi, quando eravamo matricole, avevano
abbandonato la scuola per abbracciare lo stile di vita dell'NR,
barattando carriere universitarie incerte con un'illuminazione meno
esigente.

«È solo un altro movimento millenarista,» aggiunse Jason, «in netto
ritardo per il millennio ma giusto in tempo per la fine del mondo.»

«Una setta, in altre parole.»

«No, non proprio. NR è uno slogan per tutta la gamma dei movimenti
edonisti cristiani, per cui non è una setta in sé, anche se racchiude
dei gruppi settari. Non c'è un leader unico né un testo sacro, ma solo
una banda di teologi estremisti con cui il movimento si identifica
liberamente - C.R. Ratel, Laura Greengage, persone del genere.» Avevo
visto i loro libri sugli scaffali dei negozi. La teologia dello Spin con
punti interrogativi nei titoli: \emph{Siamo stati testimoni del Secondo
	Avvento? Riusciremo a sopravvivere alla Fine dei Tempi}? «Non hanno
neanche un piano particolare, se non una specie di comunalismo da fine
settimana. Ma ciò che attira le folle non è la teologia. Hai mai visto
le riprese di questi raduni dell'NR, del tipo che loro chiamano
Ekstasis?»

Le avevo viste e, a differenza di Jase che non era mai stato molto a suo
agio con le questioni carnali, potevo comprenderne il richiamo. Avevo
visto il video di una riunione ripresa nella catena delle Cascate,
l'estate dell'anno prima. Mi era parso un incrocio tra un picnic
battista e un concerto dei Grateful Dead. Un prato assolato, fiori di
campo, tuniche bianche da cerimonia, un tizio senza un filo di grasso
che suonava uno shofar. Verso l'imbrunire un falò bruciava vivacemente
ed era stato allestito un palco per i musicisti. Poi erano cominciate a
cadere le tuniche ed era iniziata la danza. E anche qualcosa di più
intimo della danza.

A dispetto del disgusto manifestato dai media convenzionali, a me era
parso innocente e seduttivo. Niente prediche, solo qualche centinaio di
pellegrini che sorridevano alla faccia dell'estinzione e amavano il
prossimo per come voleva essere amato. Le riprese erano state
masterizzate su centinaia di DVD e passavano di mano in mano nei
dormitori delle università di tutto il Paese, compresa Stony Brook. Non
esiste atto sessuale tanto puro da impedire a uno studente single di
medicina di farsi una sega.

«È difficile immaginare che Diane possa essere attratta dall'NR.»

«Al contrario. La gente come Diane è il tipico destinatario a cui
mirano. Ha una paura matta dello Spin e di tutto ciò che implica per il
mondo. L'NR, per quelli come lei, è un calmante. Trasforma ciò che
temono di più in un oggetto di adorazione, una porta per il Regno dei
Cieli.»

«Da quanto tempo ne fa parte?»

«Ormai da quasi un anno. Da quando ha conosciuto Simon Townsend.»

«Simon è dell'NR?»

«Temo che Simon sia un \emph{integralista} dell'NR.»

«Tu l'hai conosciuto questo tizio?»

«L'ha portato alla Grande Casa lo scorso Natale. Credo che volesse
guardare i fuochi d'artificio. E.D. naturalmente non approva Simon. In
effetti, la sua ostilità era prevedibile,» qui Jason fece una smorfia al
ricordo di quella che doveva essere stata una delle più grosse sfuriate
di E.D., «ma Diane e Simon si sono comportati da bravi adepti dell'NR,
porgendo l'altra guancia. In pratica, gli hanno sorriso fin quasi ad
ammazzarlo. E lo dico in senso letterale. Un solo altro sguardo gentile
e comprensivo e lo avrebbero spedito in cardiologia per un infarto.»

``Uno a zero per Simon'' pensai. «È quello giusto per lei?»

«È proprio il tipo che cerca ma l'ultima cosa di cui ha bisogno.»

\separatorefiori

Arrivarono quel pomeriggio, scoppiettando sul viale d'accesso in un'auto
da turismo vecchia di quindici anni che sembrava bruciare più olio del
trattore di Mike-il-giardiniere. Guidava Diane. Parcheggiò e uscì
dall'altro lato della macchina, nascosta dal portapacchi, mentre Simon
era in bella vista e sorrideva timidamente.

Era un tipo di bell'aspetto. Alto un metro e ottanta o forse più, era
magro ma non gracile e con un viso comune, magari un po' equino ma
bilanciato dai capelli biondi e ribelli. Quando sorrideva, si notava una
fessura sul davanti tra i denti superiori. Indossava dei jeans, una
camicia scozzese a scacchi e una bandana azzurra legata attorno al
bicipite sinistro come un laccio emostatico; era un simbolo dell'NR,
come appresi in seguito.

Diane fece il giro dell'auto e si fermò accanto a lui, poi entrambi
salirono i gradini del portico con un ampio sorriso rivolto a me e
Jason. Anche lei era agghindata secondo l'alta moda dell'NR: una gonna
color fiordaliso lunga fino ai piedi, una camicetta azzurra e un
ridicolo cappello nero a tesa larga come quello che portano gli uomini
Amish. Ma i vestiti le stavano bene o meglio la incorniciavano in un
modo piacevole, evocando un'ottima salute e una sensualità contadina.
Aveva il viso pieno di vitalità come una bacca non colta. Si riparò gli
occhi dalla luce e sorrise, a me in particolare o almeno così volli
credere. Oddio, quel sorriso. Sincero e malizioso allo stesso tempo.

Cominciai a sentirmi perduto.

Il telefono di Jason squillò. Lo tirò fuori dalla tasca e controllò il
numero.

«A questa devo rispondere» sussurrò.

«Non lasciarmi qui da solo, Jase.»

«Vado in cucina. Torno subito.»

Se la svignò proprio quando Simon scaricò la sua sacca da viaggio sul
legno del portico dicendo: «Tu devi essere Tyler Dupree!»

Mi porse la mano. Gliela strinsi. Aveva una presa energica e un mieloso
accento del Sud, vocali trascinate come rami trasportati dalla corrente
e consonanti garbate come biglietti da visita. Pronunciò il mio nome
come se fosse cajun, ma la mia famiglia non si era mai spostata a sud di
Millinocket. Diane lo raggiunse urlando: «Tyler!» e mi strinse in un
violento abbraccio. All'improvviso mi ritrovai i suoi capelli sul viso e
riuscii a coglierne soltanto l'odore salato e solare.

Indietreggiammo lasciando tra noi la comoda distanza di un braccio.
«Tyler, Tyler,» esclamò, come se mi fossi trasformato in un tipo
eccezionale, «hai un bell'aspetto dopo tutti questi anni.»

«Otto,» dissi stupidamente, «otto anni.»

«Caspita! Davvero?»

Li aiutai a trascinare dentro i bagagli, mostrai loro il salotto
dall'ingresso e mi affrettai a rintracciare Jason che era in cucina al
cellulare. Quando entrai era girato di schiena.

«No,» disse con voce tesa, «no\ldots{} nemmeno il Dipartimento di
Stato?»

Mi fermai su due piedi. Il Dipartimento di Stato. Oddio.

«Posso rientrare fra un paio d'ore se\ldots{} ah, capisco. Va bene. No,
va bene. Ma tienimi informato. Bene. Grazie.»

Si infilò il telefono in tasca e si accorse della mia presenza.

«Stavi parlando con E.D.?» chiesi.

«Col suo assistente, in realtà.»

«Tutto bene?»

«Dai, Ty, vuoi che ti metta al corrente di \emph{tutti} i segreti?» Si
sforzò di sorridere, senza un gran esito. «Spero che tu non abbia
origliato.»

«Ho sentito solo che ti sei offerto di tornare a Washington D.C. e di
lasciarmi qui con Simon e Diane.»

«Be'\ldots{} forse dovrò. I cinesi ci stanno tirando il pacco.»

«Che intendi con \emph{tirare il pacco}?»

«Si rifiutano di abbandonare del tutto l'idea del lancio programmato.
Vogliono rinviare la decisione.»

Intendeva l'attacco nucleare agli artefatti dello Spin. «Suppongo che
qualcuno stia cercando di convincerli a non impegnarsi.»

«Sono in corso azioni diplomatiche. Però non è che stiano proprio
\emph{funzionando}. Pare che le trattative siano a un punto morto.»

«E quindi\ldots{} \emph{cazzo}, Jase! Che succede se poi lanciano
\emph{davvero}?»

«Succede che fanno detonare due armi a fusione ad alto rendimento in
prossimità dei dispositivi sconosciuti associati allo Spin. Riguardo
alle conseguenze\ldots{} be', è una domanda interessante. Ma non è
ancora successo. Probabilmente non succederà.»

«Parli della fine del mondo o forse di quella dello Spin\ldots»

«Abbassa la voce. Abbiamo ospiti, ricordi? E poi stai esagerando. Quello
che i cinesi hanno in mente è avventato e forse inutile, ma anche se
decidono di andare avanti probabilmente non sarà un suicidio. Qualunque
cosa siano gli Ipotetici, sapranno come difendersi senza dover
distruggere noi per farlo. E non è detto che gli artefatti polari siano
i dispositivi che rendono possibile lo Spin. Magari sono piattaforme di
osservazione passiva, dispositivi di comunicazione o addirittura falsi
bersagli.»

«Se i cinesi lanciano,» chiesi, «quanto preavviso abbiamo?»

«Dipende da quello che intendi con ``noi''. Il grande pubblico forse non
ne saprà nulla finché non è tutto finito.»

Fu allora che per la prima volta cominciai a capire che Jason non era
solo l'apprendista di suo padre ma aveva già iniziato a crearsi le sue
amicizie nelle alte sfere. In seguito avrei appreso molte più cose sulla
Fondazione Perielio e sul lavoro che Jason vi svolgeva. Per allora
faceva ancora parte della sua vita nell'ombra. Anche quando eravamo
bambini Jase aveva una vita nell'ombra: fuori dalla Grande Casa era un
prodigio in matematica, che frequentava senza alcuna difficoltà una
scuola privata d'elite come un campione del Masters che gioca in un
campo da minigolf; a casa era solo Jase, e ci eravamo impegnati perché
le cose restassero così.

Ed era ancora così. Solo che ora proiettava un'ombra più grande. Non
passava più le giornate a impressionare gli assistenti di calcolo alla
Rice. Ora le passava impegnandosi a influenzare il corso della storia
umana.

Aggiunse: «Se succederà, sì, mi daranno un po' di preavviso. \emph{Ci}
daranno un po' di preavviso. Non voglio però che Diane si preoccupi. E
neppure Simon, ovviamente.»

«Grandioso. Basterà smettere di pensarci. Alla fine del mondo.»

«Ma non è così, non è ancora successo nulla. Calmati, Tyler. Versa da
bere se proprio senti il bisogno di fare qualcosa.»

Per quanto si sforzasse di sembrare indifferente, gli tremava la mano
mentre prendeva quattro bicchieri dalla credenza.

Me ne sarei potuto andare. Sarei potuto uscire dalla porta, saltare
sulla mia Hyundai e allontanarmi di un bel pezzo prima che si
accorgessero della mia assenza. Pensai a Diane e Simon in salotto che
praticavano il cristianesimo hippy e a Jase in cucina che riceveva al
cellulare bollettini sulla fine del mondo: volevo davvero passare la mia
ultima notte sulla Terra con queste persone?

Eppure allo stesso tempo pensavo: ``Con chi altro sennò? Con chi
altro?''

\separatorefiori

«Ci siamo conosciuti ad Atlanta» spiegò Diane. «Lo stato della Georgia
ospitava un seminario sulla spiritualità alternativa. Simon era lì per
ascoltare la conferenza di C.R. Ratel e io l'ho incontrato per caso alla
mensa del campus. Era seduto da solo a leggere una copia del
\emph{Secondo avvento} e anch'io ero sola, perciò ho appoggiato il
vassoio sul suo tavolo e abbiamo iniziato a parlare.»

Diane e Simon erano entrambi seduti accanto alla finestra su un divano
in tessuto giallo che puzzava di polvere. Diane addossata in maniera
scomposta al bracciolo. Simon vigile e ritto. Il suo sorriso cominciava
a preoccuparmi. Non svaniva mai.

Sorseggiavamo i nostri drink mentre le tende si muovevano alla brezza e
un tafano ronzava contro la zanzariera della finestra. Era difficile
sostenere una conversazione quando c'erano troppe cose di cui non
dovevamo parlare. Mi sforzai di riprodurre il sorriso di Simon. «Quindi
studi?»

«Studiavo» puntualizzò.

«Cosa fai adesso?»

«Viaggio, per lo più.»

«Simon se lo può permettere,» disse Jase, «è un ereditiere.»

«Non essere scortese» lo ammonì Diane, e dal tono della voce trasparì un
vero e proprio avvertimento. «Per favore, Jase, almeno per stavolta.»

Ma Simon non ci fece caso. «No, diciamo che è vero. Ho dei soldi da
parte. Diane e io stiamo cogliendo l'occasione per visitare un po' il
Paese.»

«Il nonno di Simon,» annunciò Jason, «era Augustus Townsend, il re degli
scovolini per pipa della Georgia.»

Diane alzò gli occhi al cielo. Simon, ancora imperturbabile - cominciava
a sembrare quasi un santo - disse: «Un tempo era così. Ora non dobbiamo
neanche più chiamarli scovolini. Sono ``nettapipe''.» Rise. «Ed eccomi
seduto qui, l'erede di una fortuna di nettapipe.» In realtà era una
fortuna di regali e accessori, spiegò Diane in seguito. Augustus
Townsend aveva iniziato con gli scovolini, ma i soldi li aveva fatti
distribuendo giocattoli di latta, braccialetti portafortuna e pettini di
plastica per gli empori di tutto il Sud. Negli anni Quaranta la famiglia
era stata una presenza importante nei circoli sociali di Atlanta.

Jason proseguì: «Vedi, Simon ha quella che definiresti una professione.
Lui è uno spirito libero.»

«Credo che nessuno di noi sia veramente uno spirito libero,» aggiunse
Simon, «ma è vero che non ho e non desidero una professione. Immagino
che questo mi faccia apparire una persona pigra. Be', io \emph{sono}
pigro. È il mio vizio. Ma mi chiedo di quale utilità possa essere nel
lungo periodo un qualsiasi mestiere. Considerando lo stato delle cose.
Senza offesa.» Si voltò verso di me. «Sei nel campo della medicina,
Tyler?»

«Mi sono laureato da poco, ma per come vanno le cose\ldots»

«No, penso che sia meraviglioso. Forse è il lavoro più prezioso del
pianeta.»

In pratica Jason aveva accusato Simon di essere inutile. Simon aveva
risposto che le professioni in generale erano inutili\ldots{} tranne
quelle come la mia. Botta e risposta. Era come guardare una rissa da bar
con gli avversari in scarpette da ballo.

Tuttavia, sentivo di volermi scusare per Jase. Jason non era offeso
dalla filosofia di Simon ma dalla sua presenza.

Quella settimana nel Berkshire doveva essere una rimpatriata tra me,
Jason e Diane, un ritorno nella zona di comfort, una rivisitazione
dell'infanzia. Invece, ci eravamo ritrovati confinati in compagnia di
Simon, che Jason considerava un intruso, una specie di Yoko Ono con
l'accento del Sud.

Chiesi a Diane da quanto erano in viaggio.

«Da circa una settimana, ma andremo in giro per quasi tutta l'estate.
Jason ti avrà già parlato del Nuovo Regno, Ty, ma qualunque cosa ne
pensi lui, sappi che è davvero stupendo. Abbiamo una rete di amici
sparsi per tutto il Paese. Persone da cui possiamo fermarci a dormire un
giorno o due. Perciò andremo a conclavi e concerti dal Maine all'Oregon,
da luglio fino a tutto ottobre.»

Jason disse: «Immagino il risparmio su spese di alloggio e vestiti.»

«Non tutti i conclavi sono come l'Ekstasis» ribatté come un fulmine
Diane.

«Non andremo molto lontano,» aggiunse Simon, «con quella vecchia
macchina che cade a pezzi. Il motore perde colpi e finora abbiamo fatto
pochissimi chilometri. Non m'intendo molto di meccanica, purtroppo.
Tyler, capisci qualcosa di motori?»

«Un paio di cosette.» Capii che era un invito a uscire di casa con lui
mentre Diane cercava di negoziare una tregua con il fratello. «Diamo
un'occhiata.»

Fuori era ancora chiaro, ondate di aria calda si diffondevano dal prato
smeraldo oltre il vialetto. Ammetto che non ascoltai Simon con molta
attenzione mentre sollevava il cofano della sua vecchia Ford e mi
elencava i suoi problemi. Se era così ricco come aveva insinuato Jase,
perché non si comprava una macchina migliore? Mi dissi che doveva avere
già dissipato l'eredità o, magari, che la stessa fosse vincolata in
fondi fiduciari.

«So di poter sembrare stupido,» disse Simon, «soprattutto in mezzo a
voi, ma non ho mai capito molto di cose scientifiche o meccaniche.»

«Neanch'io sono un esperto. Anche se riusciamo a far andare meglio il
motore, dovresti chiedere un parere a un vero meccanico prima di guidare
per tutto il Paese.»

«Grazie, Tyler.» Mentre lo controllavo, mi osservava stralunato, quasi
in estasi. «Apprezzo il tuo consiglio.»

Era molto probabile che fosse colpa delle candele. Gli chiesi se fossero
mai state sostituite. «Non che io sappia» rispose. L'auto aveva più di
95.000 chilometri. Usai il set di chiavi a cricchetto della mia auto per
svitare una delle candele e gliela mostrai: «Ecco il problema.»

«Quell'affare?»

«E le sue compari. La buona notizia è che non sono pezzi costosi. La
brutta è che è meglio non guidarla finché non le sostituiamo.»

«Uhm» disse lui.

«Possiamo andare in città con la mia macchina e prendere i ricambi se
sei disposto ad aspettare fino a domattina.»

«Be', certo. Sei molto gentile. Non avevamo intenzione di partire
subito. Oh, a meno che Jason non insista.»

«Jason si calmerà. È solo\ldots»

«Non devi darmi spiegazioni. Jason avrebbe preferito che non fossi qui.
Lo capisco. Non mi sconvolge né mi sorprende. Diane non se la sentiva di
accettare un invito che poneva come punto fermo la mia esclusione.»

«Be'\ldots{} meglio così.» (Sì, come no.)

«Ma posso tranquillamente prenotare una stanza da qualche parte in
città.»

«Non ce n'è bisogno» dissi, chiedendomi per quale strana ragione stessi
insistendo perché Simon Townsend rimanesse. Non so cosa avessi sperato
di ottenere rincontrando Diane, ma la presenza di Simon aveva annullato
ogni speranza nascente. Forse era stato meglio così.

«Suppongo che Jason ti abbia parlato del Nuovo Regno. È stato un motivo
di scontro fra noi.»

«Mi ha detto che ne fate parte.»

«Non voglio provare a convertirti, ma se hai qualche perplessità sul
movimento forse posso rassicurarti.»

«Quello che so sull'NR è quanto vedo in televisione, Simon.»

«Alcuni lo chiamano edonismo cristiano. Io preferisco Nuovo Regno. In
poche parole il concetto è quello. Costruisci il chiliasmo vivendolo,
qui e ora. Fai in modo che l'ultima generazione sia idilliaca come la
prima.»

«Già. Be'\ldots{} Jase non ha molta pazienza con la religione.»

«No, in effetti no, ma sai una cosa, Tyler? Non credo che sia la
religione a turbarlo.»

«No?»

«No. In tutta onestà ammiro Jason Lawton, e non solo perché è
notoriamente in gamba. È un \emph{connaisseur}, se mi è concesso di
usare un parolone. Prende lo Spin sul serio. Quante persone ci sono
sulla Terra, otto miliardi? Tutti sanno che le stelle e la Luna sono
scomparse dal cielo, ma la gente continua a vivere come se questo non
fosse avvenuto. Solo pochi di noi credono davvero nello Spin. L'NR lo
prende sul serio. E così fa Jason.»

Era sconvolgente quanto questo fosse simile a quello che Jason stesso
aveva detto. «Non con lo stesso\ldots{} \emph{stile}, però.»

«È questo il punto cruciale della questione. Due visioni in competizione
per conquistarsi l'opinione pubblica. Fra non molto la gente dovrà
affrontare la realtà, volente o nolente. E dovrà scegliere fra una
comprensione scientifica e una spirituale. È questo che preoccupa Jason.
Perché quando si arriva a questioni di vita e di morte, la fede vince
sempre. Dove preferiresti trascorrere l'eternità \emph{tu}? In un
paradiso terrestre o in un laboratorio sterile?»

La risposta non mi sembrava così evidente come lo era per Simon. Mi
venne in mente la replica di Mark Twain a una domanda simile: «Il
Paradiso per il clima, l'Inferno per la compagnia.»

\separatorefiori

Dall'interno della casa provenivano i suoni di un litigio: la voce di
Diane che rimproverava Jason e le risposte astiose e rigide di suo
fratello. Io e Simon tirammo fuori dal garage un paio di sedie
pieghevoli e ci sedemmo all'ombra della tettoia per l'auto in attesa che
i gemelli finissero. Parlammo del tempo. Il tempo era molto bello. Su
quel punto raggiungemmo un accordo.

Alla fine il rumore proveniente dalla casa si placò. Dopo un po' uscì
Jason con l'aria castigata e ci invitò ad aiutarlo con il barbecue. Lo
seguimmo sul retro e mentre la griglia si scaldava continuammo a
chiacchierare piacevolmente. Diane uscì dalla casa col viso arrossato ma
l'aria trionfante. Era così che appariva ogni volta che usciva
vittoriosa da una discussione con Jase: un po' altezzosa, un po'
sorpresa.

Ci sedemmo a mangiare pollo, tè freddo e quel che era avanzato
dell'insalata ai tre fagioli. «Vi dispiace se dico una preghiera di
ringraziamento?» chiese Simon.

Jason alzò gli occhi al cielo ma annuì.

Simon chinò la testa solennemente. Mi preparai aspettandomi un sermone.
Invece disse solo: «Concedici il coraggio di accettare i doni che ci hai
posto davanti oggi e tutti gli altri giorni. Amen.»

Una preghiera che non esprimeva gratitudine ma il bisogno di coraggio.
Molto contemporanea. Diane mi sorrise dall'altra parte del tavolo. Poi
strinse il braccio di Simon e ci mettemmo a mangiare di buona lena.

\separatorefiori

Era presto quando finimmo, la luce del Sole indugiava e le zanzare
sopravvissute alla disinfestazione non avevano ancora dato il via alla
loro frenesia serale. La brezza era cessata e l'aria che si stava
rinfrescando era mite.

Altrove, le cose precipitavano.

Non sapevamo - e neppure Jason ne era a conoscenza, nonostante tutti i
suoi decantati agganci - che da qualche parte, tra quel primo boccone di
pollo e quell'ultima cucchiaiata d'insalata ai tre fagioli, i cinesi
avevano abbandonato le trattative e ordinato il lancio immediato di una
coppia di missili Dong Feng modificati, armati di testate termonucleari.
Forse i razzi stavano partendo proprio mentre prendevamo dal frigo le
Heineken. Bottiglie verdi ghiacciate a forma di razzo che grondavano
sudore estivo.

Sparecchiammo il tavolo del patio. Accennai alle candele usurate e
all'intenzione di portare Simon in città l'indomani mattina. Diane
sussurrò qualcosa a suo fratello, poi, dopo una pausa, gli diede un
colpetto con il gomito. Alla fine, Jase annuì e si rivolse a Simon: «C'è
uno di quegli ipermercati di autoricambi fuori Stockbridge che rimane
aperto fino alle nove. Se vuoi, ti ci accompagno subito.»

Era un'offerta di pace, seppur avanzata malvolentieri. Simon si riprese
dalla sorpresa iniziale e disse: «Non rifiuterò di certo un giro in
quella Ferrari, se è questa la tua offerta.»

«Posso mostrarti di cosa è capace.» Ammorbidito dalla prospettiva di
mostrare la sua auto, Jason entrò in casa per andare a prendere le
chiavi. Prima di seguirlo, Simon fece un'espressione incredula.

Guardai Diane. Sorrideva di cuore, orgogliosa di quel trionfo
diplomatico.

Altrove, i missili Dong Feng si avvicinavano e attraversavano la
barriera dello Spin in rotta verso gli obiettivi programmati. Era strano
immaginarli sfrecciare su una Terra improvvisamente buia, fredda e
immobile, in moto solo grazie a una programmazione interna, mentre
miravano a quegli artefatti anonimi sospesi alla deriva a centinaia di
chilometri sui poli.

Come un dramma senza pubblico, troppo repentino per essere visto.

\separatorefiori

Gli esperti convennero - in seguito - che la detonazione delle testate
cinesi non aveva avuto alcun effetto sul flusso differenziale del tempo.
In cambio, ciò su cui aveva inciso (e molto) era il filtro visivo che
circondava la Terra. Per non parlare della percezione umana dello Spin.

Come aveva sottolineato Jase anni prima, il gradiente temporale
significava che enormi quantità di radiazioni spostate radicalmente
verso il blu avrebbero inondato la superficie del nostro pianeta, se non
fossero state filtrate e gestite dagli Ipotetici. Oltre tre anni di luce
solare per ogni secondo trascorso: abbastanza per uccidere ogni essere
vivente sulla Terra, a rendere il terreno sterile e far ribollire gli
oceani. Gli Ipotetici, che avevano progettato l'involucro temporale
della Terra, ci avevano protetti anche dai suoi letali effetti
collaterali. Inoltre, non solo regolavano quanta energia doveva
raggiungere la Terra statica, ma anche quanta parte di calore e luce del
pianeta stesso doveva essere irradiata nello spazio. Forse era quello il
motivo per cui negli ultimi anni il tempo era stato così
piacevolmente\ldots{} regolare.

Il cielo sul Berkshire era limpido come un cristallo di Waterford quando
alle 7.55, ora locale, i missili cinesi raggiunsero i loro obiettivi.

\separatorefiori

Ero con Diane nel salotto quando squillò il telefono di casa.

Avevamo notato qualcosa prima della chiamata di Jason? Un cambiamento
nella luce, qualcosa di insignificante come la sensazione che una nuvola
potesse essere passata davanti al Sole? No, niente. Tutta la mia
attenzione era concentrata su Diane. Bevevamo drink ghiacciati e
parlavamo del più e del meno. Dei libri letti, dei film visti. La
conversazione era ipnotica, non per il contenuto ma per le cadenze del
discorso, il ritmo che adottavamo quando eravamo soli era immutato. Ogni
conversazione fra amici o amanti crea i propri ritmi distesi o
impacciati, un linguaggio celato che scorre come un fiume sotterraneo
persino sotto lo scambio più banale. Quello che ci dicevamo era trito e
convenzionale, ma il flusso invisibile era profondo e talvolta
insidioso.

E ben presto ci ritrovammo a flirtare, come se Simon Townsend e gli
ultimi otto anni non significassero niente. All'inizio per gioco, dopo
forse non più. Le dissi che mi era mancata. E lei: «Ci sono stati
momenti in cui avrei voluto parlare con te. In cui sentivo il bisogno di
parlarti. Ma non avevo il tuo numero, e immaginavo che fossi impegnato.»

«Avresti potuto trovarlo il mio numero. E non ero impegnato.»

«Hai ragione. In realtà era più\ldots{} vigliaccheria morale.»

«Faccio così tanta paura?»

«Non tu. La nostra situazione. Credo mi sentissi come se dovessi
scusarmi con te. E non sapevo da dove iniziare.» Sorrise debolmente.
«Penso di non saperlo ancora.»

«Non c'è niente di cui scusarti, Diane.»

«Grazie per averlo detto, ma non mi trovi d'accordo. Non siamo più
bambini. Ora possiamo guardarci indietro con un certo discernimento.
Eravamo vicinissimi ma in realtà non ci siamo mai toccati. Era l'unica
cosa che non potevamo fare. E neanche parlarne. Come se avessimo
prestato un giuramento di silenzio.»

«Dalla notte in cui scomparvero le stelle» dissi con la bocca secca,
sbalordito da me stesso, terrorizzato, scosso.

Diane fece un gesto con la mano. «Quella notte. Quella notte\ldots{} sai
cosa ricordo di quella notte? Il binocolo di Jason. Guardavo la Grande
Casa mentre voi due fissavate il cielo. In realtà non ricordo un bel
niente delle stelle. Invece ricordo di aver intravisto Carol in una
delle stanze da letto sul retro con qualcuno del catering. Era ubriaca e
sembrava che ci stesse provando.» Rise timidamente. «Quella fu la mia
piccola apocalisse. Tutto ciò che già odiavo della Grande Casa, della
mia famiglia, tutto si concentrò in una notte. Volevo solo fingere che
non esistesse. Nessuna Carol, nessun E.D., nessun Jason\ldots»

«Nessun me?»

Si spostò sul divano e, poiché la conversazione aveva preso una certa
piega, mi poggiò una mano sulla guancia. Aveva la mano fredda, come la
temperatura del bicchiere che aveva tenuto fino a un momento prima. «Tu
eri l'eccezione. Io ero spaventata, tu incredibilmente paziente. Lo
apprezzavo.»

«Ma non potevamo\ldots»

«Toccarci.»

«Toccarci. E.D. non l'avrebbe mai tollerato.»

Tolse la mano. «Avremmo potuto nasconderglielo, se avessimo voluto. Ma
hai ragione, il problema era E.D. Infettava tutto. Era indecente il modo
in cui faceva vivere a tua madre una specie di esistenza di seconda
classe. Era umiliante. Posso confessarlo? Odiavo essere sua figlia. E
soprattutto odiavo l'idea che, se fosse successo qualcosa tra noi,
sarebbe stato il tuo modo di vendicarti di E.D. Lawton.»

Si appoggiò allo schienale, un po' sorpresa da se stessa, credo.

«Ovviamente non sarebbe stato per quello» dissi con cautela.

«Ero confusa.»

«Per te l'NR rappresenta questo? La tua vendetta su E.D.?»

«No,» rispose continuando a sorridere, «non amo Simon solo perché fa
arrabbiare mio padre. La vita non è così semplice, Ty.»

«Non volevo insinuare che\ldots»

«Ma vedi quanto è insidiosa? Certi sospetti ti entrano in testa e
rimangono bloccati lì. No, l'NR non c'entra con mio padre. C'entra con
la volontà di trovare il divino in ciò che è successo alla Terra ed
esprimere quel divino nella vita quotidiana.»

«Forse neanche lo Spin è così semplice.»

«Simon dice che ci sta ammazzando o trasformando.»

«Mi ha detto che state costruendo il paradiso in terra.»

«Non è quello che dovrebbero fare i cristiani? Esprimere il Regno di Dio
nella vita di tutti i giorni?»

«O almeno ballare al suo ritmo.»

«Adesso mi sembri Jason. È ovvio che non possa difendere tutto quello
che succede nel movimento. La scorsa settimana eravamo a un conclave a
Filadelfia e abbiamo conosciuto una coppia, della nostra età, cordiale,
intelligente\ldots{} ``viva nello spirito'' l'ha definita Simon. Siamo
andati a cena fuori e abbiamo parlato della Parusia. Poi ci hanno
invitato nella loro camera d'albergo e all'improvviso si sono messi a
stendere strisce di cocaina e proiettare video porno. Ogni genere di
emarginato è attratto dall'NR. Senza dubbio. E per gran parte di loro la
teologia esiste a malapena, se non come immagine sfocata del Giardino
dell'Eden. Ma nella sua espressione migliore il movimento è tutto ciò
che afferma di essere, una fede viva e genuina.»

«Fede in cosa, Diane? Ekstasis? Promiscuità?»

Mi pentii di quelle parole non appena le ebbi pronunciate. Sembrava
ferita. «L'Ekstasis non ha nulla a che vedere con la promiscuità. Non
quando è vera, comunque. Ma nel corpo di Dio nessun atto è proibito,
sempre che non sia compiuto per vendetta o rabbia, e purché esprima
l'amore divino e umano.»

In quel momento squillò il telefono. Forse avevo lo sguardo colpevole.
Diane vide la mia espressione e si mise a ridere.

Quando risposi, le prime parole di Jason furono: «Ti avevo detto che
avremmo avuto un po' di preavviso. Mi dispiace, mi sbagliavo.»

«Cosa?»

«Tyler\ldots{} non hai guardato il cielo?»

\separatorefiori

Salimmo al piano di sopra per trovare una finestra che guardasse al
tramonto.

La stanza da letto a ovest era molto grande, provvista di un guardaroba
con cassettiera in mogano, un letto con le barre in ottone e grandi
finestre. Tirai le tende. Diane restò senza fiato.

Non c'era nessun tramonto. O, piuttosto, ce n'erano diversi.

L'intero cielo a ovest era illuminato. Al posto dell'unica sfera del
Sole, c'era un arco di luce rossastra che si estendeva per almeno
quindici gradi nell'orizzonte. Sembrava un'immagine tremolante a
esposizione multipla di una decina di tramonti o più. La luce era
irregolare, si ravvivava e si affievoliva come un fuoco lontano.

Restammo a guardare a bocca aperta per un tempo infinito. Alla fine
Diane disse: «Che cosa sta succedendo, Tyler? Dimmelo, che cosa sta
succedendo?»

Le raccontai quello che Jason mi aveva detto delle testate nucleari
cinesi.

«Sapeva che sarebbe potuto accadere?» mi chiese, poi si rispose da sola:
«Certo che lo sapeva.» La strana luce dava alla stanza una tonalità
rosata e le si posava sulle guance come una febbre. «Ci ucciderà?»

«Jason non crede, però farà prendere alla gente una paura del diavolo.»

«Ma c'è qualche pericolo? Di radiazioni o altro?»

Ne dubitavo, ma non era da escludere. «Accendiamo la TV» suggerii. C'era
uno schermo al plasma in ogni camera, incorniciato da pannelli in noce,
di fronte al letto. Immaginai che qualsiasi tipo di radiazione, anche
solo lontanamente letale, avrebbe distorto pure le trasmissioni e la
ricezione.

Invece, la TV funzionava abbastanza bene da mostrarci le immagini,
riprese dai canali di notizie, di folle che si radunavano nelle città di
tutta Europa, dove era già buio\ldots{} o almeno buio nella misura in
cui lo consentiva quella notte. Non ci furono radiazioni letali ma tanto
panico, destinato a crescere. Diane rimase immobile sul bordo del letto,
le mani unite sul grembo e chiaramente spaventata. Mi sedetti accanto a
lei e sussurrai: «Se avesse potuto ucciderci, a quest'ora saremmo già
morti.»

Fuori, il tramonto incespicava verso l'oscurità. Il bagliore diffuso si
dissolse in vari tramonti distinti, ognuno di un pallore spettrale, poi
in una spirale di luce solare simile a una molla luminosa che si inarcò
per tutto il cielo e svanì altrettanto improvvisamente.

Ci sedemmo fianco a fianco mentre il cielo diventava sempre più buio.

Poi spuntarono le stelle.

\separatorefiori

Riuscii a mettermi in contatto con Jase ancora una volta, prima che la
connessione crollasse. Disse che, quando il cielo si era squarciato,
Simon aveva appena pagato le candele per l'auto. Le strade fuori
Stockbridge erano già affollate e la radio riferiva di saccheggi in vari
punti di Boston e code su tutte le statali principali, perciò Jase si
era fermato in un parcheggio dietro un motel e aveva affittato una
stanza per la notte per sé e Simon. L'indomani mattina, disse,
probabilmente sarebbe dovuto tornare a Washington, ma prima avrebbe
lasciato Simon a casa.

Poi passò il suo telefono a Simon, io passai il mio a Diane e lasciai la
stanza mentre lei parlava col fidanzato. La casa estiva sembrava enorme
e vuota in modo inquietante. Andai in giro accendendo luci finché lei
non mi richiamò.

«Vuoi bere ancora?» le chiesi.

«Oh, sì» rispose.

\separatorefiori

Andammo fuori poco dopo mezzanotte.

Diane stava reagendo con coraggio. Simon le aveva fatto una specie di
pistolotto da Nuovo Regno. La teologia dell'NR non contemplava il
secondo avvento tradizionale, nessuna Assunzione in cielo o apocalisse;
lo Spin era tutte queste cose messe insieme, tutte le antiche profezie
compiute trasversalmente. E se Dio avesse voluto usare la tela del cielo
per dipingervi la nuda geometria del tempo, aveva detto Simon, lo
avrebbe fatto e la nostra soggezione e il nostro timore sarebbero stati
del tutto appropriati per l'occasione. Non dovevamo, però, farci
sopraffare da quei sentimenti, perché lo Spin era in definitiva un atto
di salvezza, l'ultimo, il miglior capitolo della storia umana.

O qualcosa di simile.

E così uscimmo a guardare il cielo perché Diane pensava che farlo fosse
coraggioso e spirituale. Il cielo era terso e l'aria odorava di pino.
L'autostrada era molto lontana, ma di tanto in tanto sentivamo fievoli
suoni di clacson e sirene.

Le nostre ombre ci danzavano intorno a seconda di quale distinta
porzione di cielo si illuminava, ora a nord, ora a sud. Ci sedemmo
sull'erba a pochi metri dal chiarore immutabile della luce del portico,
Diane si appoggiò alla mia spalla e io la cinsi con il braccio, eravamo
entrambi un po' ubriachi.

Malgrado anni di gelo emotivo, la nostra storia alla Grande Casa, il suo
fidanzamento con Simon Townsend, l'NR e l'Ekstasis e nonostante lo
sconvolgimento del cielo provocato dalle bombe atomiche cinesi, ero
squisitamente consapevole della pressione del suo corpo vicino al mio. E
la cosa strana era che mi sembrava tutto assolutamente familiare, la
curva del suo braccio sotto la mia mano e il peso della sua testa contro
la mia spalla: non era una scoperta del momento ma un ricordo. Percepivo
ogni cosa come avevo sempre immaginato che l'avrei percepita. Anche
l'odore acre della sua paura mi era familiare.

Il cielo scintillava di una strana luce. Non la luce naturale
dell'universo accelerato che ci avrebbe ucciso sul colpo. Si presentava
come una serie di istantanee del cielo, mezzenotti consecutive compresse
in microsecondi, immagini residue che svanivano come lo scoppio di una
lampadina; vedevamo prima un fotogramma del cielo, poi un altro, ma
ripreso un secolo o un millennio dopo, come sequenze di un film
surreale. Alcuni scatti sembravano lunghe esposizioni sfocate, le luci
delle stelle e il riflesso lunare apparivano come sfere spettrali, o
cerchi, o scimitarre. Altri erano fermo immagine nitidi che svanivano
rapidamente. Verso nord le linee e i cerchi nel cielo erano più stretti,
i loro raggi relativamente piccoli, mentre le stelle equatoriali erano
più inquiete e si muovevano lievi su enormi ellissi. Lune piene, mezze e
calanti brillavano intermittenti da orizzonte a orizzonte in pallide
trasparenze arancioni. La Via Lattea era un nastro di fluorescenza
bianca (ora più luminosa, ora più scura) rischiarata da stelle che si
illuminavano e si spegnevano. A ogni soffio d'aria estiva alcune stelle
nascevano e altre morivano.

E tutto si muoveva.

Si muoveva brillando al passo di un'involuta danza, che rimandava a
cicli sempre più grandi ma ancora invisibili. Sopra di noi il cielo
batteva come un cuore. «È così vivo» sospirò Diane.

C'è un pregiudizio imposto dalla nostra ridotta finestra di
consapevolezza: le cose che si muovono sono vive; quelle che non lo
fanno sono morte. Il verme vivo si attorciglia sotto la roccia morta e
statica. Le stelle e i pianeti si muovono ma solo secondo le leggi
inerti della gravitazione: una pietra può cadere ma non è viva, il moto
orbitale non è altro che la stessa caduta prolungata all'infinito.

Tuttavia, allungando la nostra breve esistenza da efemere, come avevano
fatto con noi gli Ipotetici, la differenza aveva perso nitidezza.

Le stelle nascono, vivono, muoiono e lasciano in eredità le loro ceneri
elementari alle stelle nuove. La somma dei loro moti non è semplice ma
complessa oltre i limiti dell'immaginabile, una danza di attrazione e
velocità, bella e spaventosa. Spaventosa perché le stelle che si
deformano hanno reso mutevole come un terremoto ciò che dovrebbe essere
stabile. Spaventosa perché i nostri segreti organici più profondi, gli
accoppiamenti e gli atti procreativi caotici, dopotutto, non sono così
segreti: anche le stelle sanguinano e si affaticano. ``\emph{Niente
	dura, ma tutto scorre}'', non riuscivo a ricordare dove l'avessi letto.

«Eraclito» disse Diane.

Non mi ero accorto di averlo detto ad alta voce.

«In tutti quegli anni, ai tempi della Grande Casa, in tutti quegli anni
di merda sprecati, io lo sapevo\ldots»

Le poggiai un dito sulle labbra. Potevo immaginare cosa avesse in mente.

«Voglio tornare dentro,» aggiunse, «voglio tornare in camera.»

\separatorefiori

Non chiudemmo le tendine. Le stelle rotanti e cinetiche proiettavano la
loro luce nella stanza, e nell'oscurità i disegni guizzavano sulla mia
pelle e quella di Diane come immagini sfocate, nello stesso modo in cui
le luci di una città brillano attraverso una finestra striata di
pioggia, silenziose, sinuose. Non parlammo, perché le parole sarebbero
state un impedimento. Le parole sarebbero state bugie. Facemmo l'amore
senza parole e solo alla fine mi ritrovai a pensare: ``Fa che questo
duri. Solo questo.''

Dormivamo quando il cielo si oscurò di nuovo, quando i fuochi
d'artificio celesti finalmente si affievolirono e scomparvero. L'attacco
cinese era stato poco più che un gesto. Migliaia di persone erano morte
a causa del panico globale, ma sulla Terra - e, presumibilmente, fra gli
Ipotetici - non c'erano state vittime dirette.

Il mattino dopo, il Sole spuntò in orario.

Lo squillo del telefono di casa mi svegliò. Ero solo, a letto. Diane
rispose in un'altra stanza ed entrò per avvisarmi che Jase aveva detto
che le strade erano sgombre e che stava tornando.

Diane aveva fatto la doccia, si era vestita e profumava di sapone e
cotone inamidato. «E questo è tutto?» dissi. «Simon si fa vivo e tu te
ne vai? Quello che è successo stanotte non significa niente?»

Si sedette sul letto accanto a me. «Quello che è successo stanotte non
ha mai significato che non sarei partita con Simon.»

«Pensavo significasse di più.»

«Ha significato più di quanto riesca a esprimere. Ma non cancella il
passato. Ho fatto delle promesse e ho una fede, e queste cose hanno
posto limiti ben precisi nella mia vita.»

Non sembrava convinta. Replicai: «Una fede. Non dirmi che credi a quelle
stronzate.»

Si alzò in piedi, accigliata.

«Forse no, ma forse ho bisogno di stare con qualcuno che ci crede.»

\separatorefiori

Raccolsi le mie cose e caricai il bagaglio nella Hyundai prima che Jase
e Simon tornassero. Diane mi guardava dal portico mentre chiudevo il
cofano del bagagliaio.

«Ti chiamo» promise lei.

«Fallo» le risposi.